\GalSourcePageBoundary{064}{L0063}{35}{35}

choisir une certaine racine d'un nombre entier, et regarder comme
rationnelle toute fonction rationnelle de ce radical.

Lorsque nous conviendrons de regarder ainsi comme connues
de certaines quantités, nous dirons que nous les \emph{adjoignons} à
l'équation qu'il s'agit de résoudre. Nous dirons que ces quantités
sont \emph{adjointes} à l'équation.

Cela posé, nous appellerons \emph{rationnelle} toute quantité qui
s'exprimera en fonction rationnelle des coefficients de l'équation
et d'un certain nombre de quantités \emph{adjointes} à l'équation et
convenues arbitrairement.

Quand nous nous servirons d'équations auxiliaires, elles seront
rationnelles, si leurs coefficients sont rationnels en notre sens.

On voit, au surplus, que les propriétés et les difficultés d'une
équation peuvent être tout à fait différentes suivant les quantités
qui lui sont adjointes. Par exemple, l'adjonction d'une quantité
peut rendre réductible une équation irréductible.

Ainsi, quand on adjoint à l'équation
\[
\frac{x^{n}-1}{x-1}=0\GalControlReading{.}{W09-CD007}
\]
où $n$ est premier, une racine d'une des équations auxiliaires de
M.~Gauss, cette équation se décompose en facteurs et devient, par
conséquent, réductible.

Les substitutions sont le passage d'une permutation à l'autre.

La permutation d'où l'on part pour indiquer les substitutions
est toute arbitraire, quand il s'agit de fonctions; car il n'y a aucune
raison pour que, dans une fonction de plusieurs lettres, une
lettre occupe un rang plutôt qu'un autre.

Cependant, comme on ne peut guère se former l'idée d'une
substitution sans se former celle d'une permutation, nous ferons,
dans le langage, un emploi fréquent des permutations, et nous ne
considérerons les substitutions que comme le passage d'une permutation
à une autre.

Quand nous voudrons grouper des substitutions, nous les ferons
toutes provenir d'une même permutation.

Comme il s'agit toujours de questions où la disposition primitive
des lettres n'influe en rien dans les groupes que nous considérerons,
on devra avoir les mêmes substitutions, quelle que soit
